\section{Add the original state and serialize the result}
The working words after twenty rounds are not yet the output block. Add each
original word to its corresponding working word, retaining the low 32 bits, and
write that sum in little-endian byte order. This final addition is called
\emph{feed-forward}. Its inputs are the original state and the fully transformed
working state. Adding the original words once per double-round would define a
different calculation.

In the scalar version, \texttt{s0} points to the preserved input state and
\texttt{s1} points to the working output. With a byte offset in \texttt{t0}, one
iteration loads the matching words, adds them, and overwrites the working word.
\Needspace{11\baselineskip}
\MARSCodeLabel{Publish one feed-forward word on the little-endian target.}
\begin{minted}[samepage=true]{mars-idea}
daddu $t1, $s0, $t0
lw $t2, 0($t1)
daddu $t1, $s1, $t0
lw $t3, 0($t1)
addu $t2, $t2, $t3
sw $t2, 0($t1)
\end{minted}

Advance the offset by four and stop at sixty-four. The input/output separation
allows every original word to remain intact until it is used. Bytes beyond the
64-byte output are outside the function's ownership, including any convenient
padding allocated by a caller.

For the block example in RFC 8439, section 2.3.2, use key bytes zero through 31,
counter one, and nonce bytes shown below. The nonce differs from the RFC's later
whole-message cipher example; copying that example's nonce would change the block.
\begin{verbatim}
00 00 00 09 00 00 00 4a 00 00 00 00
\end{verbatim}
The first working word after twenty rounds is \texttt{837778ab}. Adding its original
constant \texttt{61707865} gives \texttt{e4e7f110} after the carry is discarded.
The first four output bytes are therefore \texttt{10 f1 e7 e4}. The following word
is \texttt{15593bd1}, whose bytes are \texttt{d1 3b 59 15}. A word display and a
byte display show the same storage in different forms. A comparison that reverses
only the displayed hexadecimal string can hide an actual serialization error;
compare the bytes that the XOR routine will load.
